\documentclass{article}
\usepackage{amsmath}
\usepackage{amsthm}
\usepackage{amssymb}
\usepackage{pgfplots}
\usepackage{amsfonts}
\usepackage[margin=2.5cm]{geometry}
\usepackage{graphicx}
\usepackage[export]{adjustbox}
\usepackage{fancyhdr}
\usepackage[portuguese]{babel}
\usepackage{hyperref}
\usepackage{lastpage}
\usepackage{mathtools}
\usepackage{subfiles}

\pagestyle{fancy}
\fancyhf{}

\pgfplotsset{compat = 1.18}

\hypersetup{
    colorlinks,
    citecolor=black,
    filecolor=black,
    linkcolor=black,
    urlcolor=black
}

\newtheorem*{def*}{\underline{Defini\c c\~ao}}
\newtheorem*{theorem*}{\underline{Teorema:}}
\newtheorem{example}{\underline{Exemplo:}}[section]
\newtheorem*{proposition*}{\underline{Proposi\c c\~ao:}}
\newtheorem*{corol*}{\underline{Corol\'ario:}}
\newtheorem*{sol*}{\underline{Solu\c c\~ao:}}
\newtheorem*{proof*}{\underline{Prova:}}
\renewcommand\qedsymbol{$\blacksquare$}
\newcommand{\Lin}[1]{Lin_{\mathbb{K}}({#1})}

\rfoot{P\'agina \thepage \hspace{1pt} de \pageref{LastPage}}

\title{EXERC\'ICIOS DE C\'ALCULO}
\author{Renan Wenzel}
\date{\today}

\begin{document}

\begin{figure}[ht]
	\minipage{0.76\textwidth}
	\includegraphics[width=4cm]{icmc.png}
	\hspace{5cm}
	\includegraphics[height=4.9cm,width=4cm]{brasao_usp_cor.jpg}
	\endminipage
\end{figure}

\begin{center}
	\vspace{1cm}
	\LARGE
	UNIVERSIDADE DE S\~AO PAULO

	\vspace{1.3cm}
	\LARGE
	INSTITUTO DE CI\^ENCIAS MATEM\'ATICAS E COMPUTACIONAIS - ICMC

	\vspace{1.7cm}
	\Large
	\textbf{ANOTA\c C\~OES E RESUMOS DE \'ALGEBRA LINEAR}

	\vspace{1.3cm}
	\large
	\textbf{Renan Wenzel - 11169472}

	\vspace{1.3cm}
	\today
\end{center}

\newpage

\tableofcontents

\newpage

\section{Introdu\c c\~ao}

A \'algebra linear se preocupa em generalizar alguns conceitos fundamentais da matem\'atica.
Durante o ensino m\'edio e no primeiro semestre, o estudante aprendeu sobre a estrutura
dos n\'umeros inteiros, reais, racionais e, talvez, sobre os complexos. Agora, vamos ver o que
tem de mais profundo por tr\'as deles, estudando, por exemplo, quest\~oes da dimens\~ao, de
opera\c c\~oes e dos objetos que sofrem delas.

Na segunda parte do curso, o estudante costuma aprender sobre funcionais e transforma\c c\~oes
lineares, que relacionam dois espa\c cos vetoriais (ou um corpo) e satisfazem algumas condi\c c\~oes.
Um exemplo j\'a visto de uma transforma\c c\~ao linear \'e a derivada, que serve de motiva\c c\~ao
para uma constru\c c\~ao futura chamada Espa\c co Tangente.

Saber manipular e trabalhar com esses conceitos \'e essencial na carreira de matem\'atico, porque
muitos conceitos futuros tomam como base o que ser\'a visto aqui, em todas as \'areas da matem\'atica e,
para isso, espero que essas notas sirvam como um guia para quem necessita de uma luz nesse curso.
\newpage

\subfile{classes/vector_spaces}
\newpage
\subfile{classes/linear_transformations}
\newpage

\section{Polin\^omios Caracter\'isticos}

\end{document}
